% =====================================================================
% Da inserire in §V-D, dopo il paragrafo sulla calibrazione.
% Numeri: 30 coppie sorgente-target, i3lung_*_shiftdiag.csv e
% i3lung_*_summary.csv, rerun del 2026-08-21. Spearman, descrittivi.
% =====================================================================

Which transfers will arrive miscalibrated is visible in one diagnostic and
invisible in the other. Across the thirty pairs, the gap in outcome prevalence
between source and target tracks ERM's calibration error on the target
($\rho=0.50$), while the covariate distance between the two cohorts carries no
association with it ($\rho=-0.15$). The two diagnostics are themselves unrelated
here ($\rho=-0.01$): centers with nearly identical covariate distributions
differ widely in outcome rate, and the reverse holds as well. We report these
associations descriptively, since the thirty transfers share six centers and are
not independent.

The prevalence gap is a symptom and we do not read a mechanism into it. A
difference in outcome rate can follow from treatment practice, from how the
outcome is recorded, or from a genuine difference in baseline hazard, and the
three are not separable from the quantity itself. What matters for deployment is
that computing it requires the target labels, which is the information a team
planning a transfer does not have. The diagnostic that is available before
deployment is the covariate distance, and it predicts neither the discrimination
a model will lose \mbox{($\rho=-0.41$)} nor the calibration error it will carry.

The same gap bounds what source-fitted recalibration can do. Platt scaling
estimates its correction on source data, so a target whose outcome rate differs
receives a correction fitted to the wrong prevalence. Its gain falls as the gap
widens ($\rho=-0.53$), which places the same boundary the synthetic benchmark
draws when prior adjustment is the only method that improves calibration under
label-prior shift and degrades it everywhere else.
